\subsection{SQUARE MATRICES}

DEFINITION 4. A matrix whose rows and columns have the same indexing set is called a square matrix.

A square matrix with n rows and n columns is called a matrix of order n.

Remark. It should be noted that a matrix for which the indexing sets of the rows and columns have the same cardinal but are not identical, must not be considered as a square matrix; in particular, the product of two such matrices over a ring is not defined.

Clearly addition and multiplication of square matrices over A with a finite set as indexing set of the rows and columns, define on the set of these matrices a ring structure because of formulae (7), (8) and (9) (no. 3); the matrix \((\delta_{ij})\) , where \(\delta_{ij}\) is the Kronecker index (for \(i \in I\) , \(j \in I\) ), is the unit element of this ring and is denoted by \(I_n\) or \(1_n\) when I has \(n\) elements. When \(\mathbf{I} = [1, n]\) , the ring of matrices thus defined is denoted simply by \(\mathbf{M}_n(\mathbf{A})\) ; the group of invertible elements of \(\mathbf{M}_n(\mathbf{A})\) is denoted by \(\mathbf{GL}_n(\mathbf{A})\) or \(\mathbf{GL}(n, \mathbf{A})\) .

For a square matrix \(U = (a_{ij})\) of order n over A to be right (resp. left) invertible, it is necessary and sufficient that, for every system \((b_i)_{1 \leq i \leq n}\) of elements of A, the system of n equations in n unknowns

\[
\sum_ {j = 1} ^ {n} a _ {i j} x _ {j} = b _ {i} \quad (1 \leqslant i \leqslant n)
\]

\[
\left(\text {resp.} \sum_ {j = 1} ^ {n} x _ {j} a _ {j i} = b _ {i}\right)
\]

have one solution \((x_{i})\) in A.

Let I be a finite indexing set, A a ring and E a right (resp. left) A-module with basis \((e_{i})_{i\in I}\) . For every endomorphism u of E, the matrix \(M(u)\) of u with respect to the two bases identical with \((e_{i})\) is a square matrix; more briefly, it is called the matrix of u with respect to the basis \((e_{i})\) .

Suppose that \(\mathbf{I} = [1, n]\) . The mapping \(u \mapsto M(u)\) (resp. \(u \mapsto {}^t M(u)\) ) is an isomorphism of the ring \(\operatorname{End}_{\mathbf{A}}(\mathbf{E})\) onto \(\mathbf{M}_n(\mathbf{A})\) (resp. onto the opposite ring of \(\mathbf{M}_n(\mathbf{A})\) , as follows from formulae (18 D) (resp. 18 G)) (no. 4). The invertible elements of the ring \(\mathbf{M}_{n}(\mathbf{A})\) called invertible matrices, correspond under the mapping \(u \mapsto M(u)\) (resp. \(u \mapsto {}^{t}M(u)\) ) to the automorphisms of E; the group \(\mathbf{GL}(n, \mathbf{A})\) is therefore canonically identified with the group \(\mathbf{GL}(\mathbf{A}_{d}^{n})\) .

If \(u\) is an automorphism of \(\mathbf{E}\) , its contragredient \(\check{u}\) is an automorphism of the left (resp. right) A-module \(\mathbf{E}^*\) , such that \(\check{u} = (^{t}u)^{-1} = ^{t}(u^{-1})\) (§ 2, no. 5, Definition 6); if \(M(\check{u})\) is the matrix of \(\check{u}\) with respect to the dual basis \((e_{i}^{*})\) , then, by virtue of Proposition 3 (no. 4),

\[
M (\check {u}) = (^ {t} M (u)) ^ {- 1} = ^ {t} M (u ^ {- 1}).\tag{29}
\]

For every invertible matrix X, it therefore follows that \({}^{t}(X^{-1}) = ({}^{t}X)^{-1}\) ; this matrix is also denoted by \({}^{t}X^{-1}\) and called the contragredient of the matrix X.

Let \(\sigma\) be an automorphism of the ring A; for every semi-linear mapping \(u: \mathbf{E} \to \mathbf{E}\) relative to \(\sigma\) , the matrix \(M(u)\) of this mapping with respect to a basis \((e_i)\) of E is also a square matrix. It follows immediately from (27 D) (no. 6) that, if \(u\) is bijective, then

\[
M (u ^ {- 1}) = (\sigma^ {- 1} (M (u))) ^ {- 1}.
\]

Let E be an A-module which is the direct sum of a finite family \((\mathbf{E}_{i})_{i\in I}\) of submodules; for every endomorphism u of E, the matrix \(M(u)=(u_{kt})\) of u with respect to the two decompositions of E identical with \((\mathbf{E}_{i})\) (no. 5) is a square matrix of linear mappings. In order that \(u(\mathbf{E}_{i})\subset\mathbf{E}_{i}\) for all \(i\in I\) , it is necessary and sufficient that \(u_{kt}=0\) for \(k\neq i\) . When \(\mathbf{I} = [1, n]\) , the relations

\[
u (\mathrm{E} _ {i}) \subset \mathrm{E} _ {i} + \mathrm{E} _ {i + 1} + \dots + \mathrm{E} _ {n} \quad (1 \leqslant i \leqslant n)
\]

are equivalent to the relations \(u_{ki} = 0\) for \(k < i\) .

Examples of square matrices. I. Diagonal matrices. In a square matrix

\[
M = \left(m _ {\iota \kappa}\right) _ {(\iota, \kappa) \in I \times I},
\]

the elements both of whose indices are equal are called diagonal elements and the family \((m_{\iota})_{\iota \in I}\) is called the diagonal of \(M\) ; a square matrix \(M = (m_{\iota \kappa})\) over a ring, whose elements other than the diagonal elements are zero, is called a diagonal matrix. For every family \((a_{\iota})_{\iota \in I}\) of elements of a ring A, the diagonal matrix \((m_{\iota \kappa})\) such that \(m_{\iota} = a_{\iota}\) for all \(\iota \in I\) is denoted by \(\operatorname{diag}(a_{\iota})_{\iota \in I}\) (or \(\operatorname{diag}(a_1, a_2, \ldots, a_n)\) when \(I = [1, n]\) ). In the set \(\mathbf{M}_n(A)\) of square matrices of order \(n\) over A, the unit matrix \(I_n\) is a diagonal matrix and also every multiple \(aI_n = I_na\) of this matrix by a scalar \(a\) (the diagonal matrix (called scalar) all of whose diagonal elements are equal to \(a\) ).

For every family \((d_i)_{1\leqslant i\leqslant n}\) of elements of A and every matrix \(X = (x_{ij})\) of type \((n,q)\) (resp. \((p,n))\) over A, writing \(D = \mathrm{diag}(d_i)\) ,

\[
\left\{ \begin{array}{l} D X = (d _ {i} x _ {i j}) \\ X D = (x _ {i j} d _ {j}). \end{array} \right.\tag{30}
\]

In particular, for two diagonal matrices of order n,

\[
\begin{array}{c} \operatorname{diag} (a _ {i}) + \operatorname{diag} (b _ {i}) = \operatorname{diag} (a _ {i} + b _ {i}) \\ \operatorname{diag} (a _ {i}). \operatorname{diag} (b _ {i}) = \operatorname{diag} (a _ {i} b _ {i}). \end{array}\tag{31}
\]

The diagonal matrices therefore form a subring of \(\mathbf{M}_{n}(\mathbf{A})\) isomorphic to the product ring \(A^{n}\) ; the scalar matrices form a subring isomorphic to A.

\begin{enumerate}
\def\labelenumi{\Roman{enumi}.}
\setcounter{enumi}{1}
\tightlist
\item
  Permutation matrices; monomial matrices. Let \(\pi\) be any permutation of a finite set I and let \((e_i)_{i\in I}\) be the canonical basis of the A-module \(\mathbf{E} = \mathbf{A}_d^{\mathbf{I}}\) ; there exists one and only one endomorphism \(u_{\pi}\) of E such that, for all \(i\in I\) , \(u_{\pi}(e_i) = e_{\pi(i)}\) (§ 1, no. 11, Corollary 3 to Proposition 17). For all \(i\in I\) , the column of index \(i\) in the matrix \(M(u_{\pi})\) with respect to the basis \((e_i)\) has all its elements zero except the one in the row of index \(\pi(i)\) , which is equal to 1. By an abuse of language, \(M(u_{\pi})\) is called the matrix of the permutation \(\pi\) . It is immediate that for any two permutations \(\sigma, \tau\) of I, \(u_{\sigma \tau} = u_{\sigma} \circ u_{\tau}\) and that for the identity permutation \(\varepsilon, u_{\varepsilon}\) is the identity; the mapping \(\pi \mapsto M(u_{\pi})\) is therefore an isomorphism of the symmetric group \(\mathfrak{S}_{\mathbf{I}}\) onto the group of permutation matrices.
\end{enumerate}

Each row and each column of a permutation matrix contains only a single element \(\neq0\) . A finite square matrix R over a non-zero ring A, with this property, is called a monomial matrix; let \(r_{i}\) be the unique element \(\neq0\) in the column of R of index i and let \(\pi(i)\) be the index of the row where this element is; clearly \(\pi\) is a permutation of the indexing set I and \(R=M(u_{\pi})D\) , where \(D=\mathrm{diag}(r_{i})\) .

\begin{enumerate}
\def\labelenumi{\Roman{enumi}.}
\setcounter{enumi}{2}
\tightlist
\item
  Triangular matrices. In the ring \(\mathbf{M}_{n}(\mathbf{A})\) of square matrices of order n over a ring A, any matrix \((a_{ij})\) such that \(a_{ij}=0\) for i\textgreater j (resp. i\textless j) is called an upper (resp. lower) triangular matrix; it is also said that such a matrix has only zeros below (resp. above) its diagonal. It is immediately established that the upper (resp. lower) triangular matrices form a subring S (resp. T) of \(\mathbf{M}_{n}(\mathbf{A}),\mathbf{S}\cap\mathbf{T}\) being obviously the ring of diagonal matrices.
\end{enumerate}

The set S' (resp. T') of matrices in S (resp. T) whose diagonal elements are invertible is a multiplicative group of matrices called the upper (resp. lower) total triangular group, this follows immediately from § 1, no. 11, Remark 5. The set \(S_{1}\) (resp. \(T_{1}\)) of matrices in S (resp. T) whose diagonal elements are all equal to 1 is a subgroup of the above group, called the upper (resp. lower) strict triangular group, and every matrix M \ensuremath{\in} S' (resp. M \ensuremath{\in} T') whose diagonal is \((d_{i})\) , may be written as \(M = DM_{1} = M_{1}'D\) , where \(D = \mathrm{diag}(d_{i})\) and \(M_{1}\) and \(M_{1}'\) matrices belonging to \(S_{1}\) (resp. \(T_{1}\)).

\begin{enumerate}
\def\labelenumi{\Roman{enumi}.}
\setcounter{enumi}{3}
\tightlist
\item
  Diagonal and triangular matrices of matrices. Let \((\mathbf{I}_k)_{1 \leqslant k \leqslant p}\) be a partition of the finite set I; every square matrix over a ring A with indexing set I can be written in the form of a square matrix of matrices corresponding to the same partition \((I_{k})\) of the indexing set of the rows and the indexing set of the columns (no. 5)
\end{enumerate}

\[
\left( \begin{array}{c c c c} X _ {1 1} & X _ {1 2} & \ldots & X _ {1 p} \\ X _ {2 1} & X _ {2 2} & \ldots & X _ {2 p} \\ \cdot & \cdot & \cdot & \cdot \\ X _ {p 1} & X _ {p 2} & \ldots & X _ {p p} \end{array} \right)\tag{32}
\]

where each \(X_{kk}\) is a square matrix with \(I_{k}\) as indexing set of the rows and columns.

With this notation, (32) will be called a diagonal (resp. upper triangular, resp. lower triangular) matrix of matrices if all the matrices \(X_{ij}\) such that \(i \neq j\) (resp. i \textgreater{} j, resp. i \textless{} j) are zero. The interpretation of endomorphisms u whose matrix is a diagonal, resp. triangular, matrix of matrices has been seen earlier, by considering the corresponding matrix \(M(u)\) of linear mappings. The lower triangular (resp. upper triangular, diagonal) matrices of matrices for a given partition \((\mathbf{I}_{k})\) of I form subrings of the ring of matrices \(A^{I \times I}\) . In particular, the ring of diagonal matrices of matrices relative to

the partition \((\mathbf{I}_k)\) is isomorphic to the product \(\prod_{k=1}^{p} \operatorname{End}_{\mathbf{A}}(\mathbf{E}_k)\) .
% semantic-fragments: legacy-input-seam
\relax{}
